\documentclass[11pt,letterpaper]{article} % Page layout
\usepackage[margin=0.8in]{geometry}
\usepackage{setspace, indentfirst}
\usepackage{amsmath, amssymb, amsfonts, amsthm, bm, commath} % Math
\usepackage{graphicx, subcaption, wrapfig, caption} % Figures
\usepackage{sidecap, caption}
\usepackage{xcolor, enumitem} % Text / formatting
\usepackage[small,compact]{titlesec}
\usepackage[superscript,biblabel]{cite}
\usepackage{hyperref, url} % Citations / links
\newcommand{\inlinecomment}[1]{}

\pagenumbering{gobble}
\topmargin=-0.45in
\headsep=0.25in
\title{\vspace{-2.5cm} WASABI (Input output / one timepoint version) \\
  \large  \vspace{-0.4cm}}
% \author{Student: Bryan Tseng, Advisors: Adam S. Charles}
\date{\vspace{-0.5cm}}
\begin{document}

\maketitle
\vspace{-0.7cm}
\subsection*{Simplified formulation}
\begin{align*}
\mathbf{r} &\sim \mathrm{Poisson}(\boldsymbol{\lambda}) \\
\boldsymbol{\lambda}^{\mathrm{stim}} &= \mathbf{Cs} + \boldsymbol{\lambda}^{\mathrm{control}} \\
\mathbf{C} &= \mathbf{DW}_{\text{low-rank}}
\end{align*}
\begin{align*}
\mathbf{r} &:\ \text{Recorded responses (spikes)}\\
\boldsymbol{\lambda}  &:\ \text{Firing rates} \\
\mathbf{C} &:\ \text{Mapping from stimulation channels to firing rates} \\
\mathbf{s} &:\ \text{Selection of stimulation channels}\\
\mathbf{D} &:\ \text{Mapping from latent variables to recorded responses} \\
\mathbf{W}_{\text{low-rank}} &:\ \text{Low rank mapping from stimulation channels to latent variables}
\end{align*}

\begin{itemize}
\item $\mathbf{W}$ is unknown.
\item At time t-1, we want to find $\mathbf{s}_t$ that maximizes the information about $\mathbf{W}$, given $\mathbf{r}_{1:t-1}$ and $\mathbf{s}_{1:t-1}$
\end{itemize}

\clearpage
\subsection*{Ridge regression for learning W at time t}
If we assume we know $D_t$ at time $t$, we estimate $\mathbf{W}_t$ using a ridge least-squares objective:
\begin{align*}
\hat{\mathbf{W}}_t &= \arg\min_{\mathbf{W}}
||\mathbf{Z-WS}||_F^2 + \beta||\mathbf{W}||_F^2
\end{align*}
where $\beta$ is the weight of the ridge prior. We can expanding the Frobenius norm using trace:
\begin{align*}
\mathcal{L}(\hat{\mathbf{W}}_t) &= tr[(\mathbf{Z-WS})(\mathbf{Z-WS})^{\top}]+\beta \,tr(\mathbf{WW}^{\top})\\
\mathcal{L}(\hat{\mathbf{W}}_t)&=tr(\mathbf{ZZ}^{\top})
- 2\,tr(\mathbf{ZS^{\top}W^{\top}}) + tr(\mathbf{WSS^{\top}W^{\top}})
+ \beta \,tr(\mathbf{WW}^{\top})
\end{align*}
Keep terms relative to $\mathbf{W}_t$
\begin{align*}
\mathcal{L}(\hat{\mathbf{W}}_t)= -2\,\mathrm{tr}(\mathbf{ZS}^{\top}\mathbf{W}^{\top})
+ \mathrm{tr}(\mathbf{WSS^{\top}W^{\top}})
+ \beta \,\mathrm{tr}(\mathbf{WW}^{\top})
\end{align*}
Taking derivative and set to zero:

\begin{align*}
\frac{\partial \mathcal{L}}{\partial \mathbf{W}} &= -2\,\mathbf{ZS}^{\top}
+ 2\,\mathbf{WSS}^{\top}
+ 2\beta\,\mathbf{W}\\
0 &= -2\,\mathbf{ZS}^{\top} + 2\,\mathbf{WSS}^{\top}+ 2\beta\,\mathbf{W}\\
\mathbf{WSS}^{\top} &+ \beta \mathbf{W} = \mathbf{ZS}^{\top}\\
\mathbf{W}(\mathbf{SS}^{\top} &+ \beta \mathbf{I}) = \mathbf{ZS}^{\top}
\end{align*}
\begin{align*}
\boxed{\mathbf{W}=\mathbf{ZS}^{\top}(\mathbf{SS}^{\top}+\beta \mathbf{I})^{-1}}
\end{align*}

\subsection*{Generative model for perturbation}
\begin{align*}
(\boldsymbol{\lambda}_{\text{stim}}-\boldsymbol{\lambda}_{\text{control}})=\mathbf{D}(\mathbf{WS}+\boldsymbol{\xi})+\boldsymbol{\eta}
\end{align*}
Note:
\begin{itemize}
\item Current loss is written for estimating $\mathbf{W}_t$
\item $\mathbf{D}$ - learned from control? or learn jointly with trials t (Alternateing or SVD)?
\item $\boldsymbol{\xi}_t$ is perturbation noise (e.g.\ mismatch between intended and actual stim-response effects) - hyperparameter, match data?
\item $\boldsymbol{\eta}$: observation noise - hyperparameter, match data?
\item Columns of $\mathbf{W}$ have unit variance
\item Ridge prior should scale as $\beta \propto \frac{1}{n}$ so in low-trial regime solution is not degenerate and high trial regime data info is respected.
\end{itemize}
\clearpage
\subsection*{Minimizing the posterior covariance of W (Gaussian)}
For adaptive selection of stimulation patterns, at each trial $t$, we want to select $\mathbf{s}_{t}$ that maximizes the expected information gain about $\mathbf{W}$. Equivalently, this would be selecting $\mathbf{s}_{t}$ that minimizes the determinant of the covariance of $\mathbf{W}$ at time t. We separate out the Poisson MLE rate estimation and identification of the optimal $\mathbf{W}^*$ in our formulation. This formulation could be good for control-theoretic implementations later too.

\begin{equation*}
    \boldsymbol{\lambda}_t =\mathbf{DWs}_t+\boldsymbol{\lambda}^{control}_{t}
\end{equation*}

Here, we are trying to find the optimal $\mathbf{W}^*$, where $\boldsymbol{\lambda}$, $\mathbf{s}_t$, and $\boldsymbol{\lambda}^{control}_{t}$ are known.
\begin{align*}
    \text{Let } & \Sigma_{t-1}
    && \text{be the covariance matrix of $\mathbf{W}$ before trial $t-1$}\\
    \text{Let } & \beta
    && \text{be how we scale our data relative to our prior (could be $1/n$)}
\end{align*}

In terms of precision, we can formulate the Gaussian case as:
\begin{equation*}
    \Sigma_{t}^{-1} = \Sigma_{t-1}^{-1} + \beta \; \mathbb{E}[\boldsymbol{\lambda} _t\mid \mathbf{s}_{t}] \mathbb{E}[\boldsymbol{\lambda} _t\mid \mathbf{s}_{t}]^T
\end{equation*}

Take the determinant
\[
    \text{det}\;\Sigma_{t}^{-1} = \text{det}(\Sigma_{t-1}^{-1} + \beta \; \mathbb{E}[\boldsymbol{\lambda} _t\mid \mathbf{s}_{t}] \mathbb{E}[\boldsymbol{\lambda} _t\mid \mathbf{s}_{t}]^T)
\]

using the matrix inversion lemma:
\[
    \text{det}(\mathbf{A} + \mathbf{uv}^T) = (1 + \mathbf{v}^T\mathbf{A}^{-1}\mathbf{u})\text{det}(\mathbf{A})
\]
\[
    \text{det}(\Sigma_{t-1}^{-1} + \beta \; \mathbb{E}[\boldsymbol{\lambda} _t\mid \mathbf{s}_{t}] \mathbb{E}[\boldsymbol{\lambda} _t\mid \mathbf{s}_{t}]^T) =  (1+ \beta \; \mathbb{E}[\boldsymbol{\lambda} _t\mid \mathbf{s}_{t}]^T \Sigma_{t-1} \mathbb{E}[\boldsymbol{\lambda} _t\mid \mathbf{s}_{t}])\det(\Sigma_{t-1}^{-1})
\]

Since we minimize over $\mathbf{s}_{t}$, the objective function is therefore:
\begin{align*}
    \arg\max_{\mathbf{s}}\; \det\big(\Sigma_{t}^{-1}\big)
    &= \arg\max_{\mathbf{s}}\;  \big(1+ \beta \,\mathbb{E}[\boldsymbol{\lambda} _t\mid \mathbf{s}_{t}]^{\top}\Sigma_{t-1}\mathbb{E}[\boldsymbol{\lambda} _t\mid \mathbf{s}_{t}]\big)\,\det(\Sigma_{t-1}^{-1})\\
    &= \arg\max_{\mathbf{s}}\; \beta \,\mathbb{E}[\boldsymbol{\lambda} _t\mid \mathbf{s}_{t}]^{\top}\Sigma_{t-1}\mathbb{E}[\boldsymbol{\lambda} _t\mid \mathbf{s}_{t}]
\end{align*}

We can further expand out:
\begin{align*}
    \mathbb{E}[\boldsymbol{\lambda} _t\mid \mathbf{s}_{t}]=\hat{\mathbf{D}}_{t-1}\hat{\mathbf{W}}_{t-1}\mathbf{s}_{t}+\boldsymbol{\lambda}_{t}^{control}
\end{align*}
where $\hat{\mathbf{D}}_{t-1}$ and $\hat{\mathbf{W}}_{t-1}$ are estimated from timepoints $i = 1,\ldots, t-1$
and $\boldsymbol{\lambda}_{t}^{control}$ is firing rate identified from control reaches.

Need to find the correct way to express how s affects the uncertainty of W
Then take the nuclear norm and then max min (E optimality)
max min singular value of precision 

\vspace{-1cm}
\end{document}